\section{Introduction}
\label{sec:intro}

The next decade will place several independent positioning, navigation, and timing
(PNT) services around the Moon. NASA's LunaNet and Lunar Communications Relay and
Navigation Systems, ESA's Moonlight programme and its Lunar Communications and
Navigation Service (LCNS), and Japan's Lunar Navigation Satellite System are each
designed to broadcast navigation signals to surface and orbital users
\citep{israel2020lunanet,fienga2024moonlight}. The stated intent is
\emph{interoperability}: a user should be able to consume signals from more than one
provider, or hand a position off between providers, and obtain a single consistent
fix. Realizing that intent requires two distinct kinds of agreement. Providers must
agree on \emph{how} they talk---the message formats, the frame and time tags, the
coordinate conventions---and they must agree on \emph{what} they broadcast, closely
enough that a user fusing them does not accumulate error from their disagreement.

The interoperability standards now being written address the first kind of agreement
in detail and the second kind not at all. The LunaNet Interoperability
Specification (LNIS) and the supporting CCSDS/IOAG cislunar work
\citep{lnis2023,ioag2022lunar} define lunar reference systems, frame realizations,
timescales, and message formats; they make providers mutually \emph{intelligible}.
What they do not yet publish is a cross-provider \emph{error budget}: the maximum
amount by which two compliant providers may disagree---in frame realization and in
broadcast dynamics---before a multi-provider user is pushed outside its accuracy
requirement. In the applicable-document structure of the specification this quantity
is deferred to AD-5, the Lunar Reference System and LunaNet Reference Time System
Standard ({LNIS-TBD-AD0005}), which is not yet released and does not yet specify a
cross-provider consistency tolerance. Tellingly, the scope AD-5 does declare---a
lunar gravity field to finite degree and order together with a parametric
precession/nutation model---is precisely the lunar-orbit dynamics that this paper
identifies as the irreducible cross-provider floor. Supplying the number is hard
because \emph{no two lunar-PNT providers have flown}; there is, as yet, no measured
cross-provider disagreement to budget against.

Our route around that obstacle is to observe that the ``multiple providers'' problem
already exists in a real, measurable form one level down. Three independent
institutions maintain authoritative, continuously updated lunar and planetary
ephemerides using independent data reductions and dynamical models: DE440 from
JPL/NASA \citep{park2021de440}, INPOP21a from IMCCE/Observatoire de Paris
\citep{fienga2021inpop21a}, and EPM2021 from IAA~RAS \citep{pitjeva2021epm2021}. Each
realizes a lunar reference frame and a lunar orbit; each disagrees with the others by
a real, decomposable amount. That disagreement is the closest existing analogue of a
cross-provider inconsistency floor, and it can be measured today with no flight data
whatsoever. We use it to anchor the budget, while stating plainly (throughout) that
ephemerides are a representative analog for providers, not the providers themselves.

\paragraph{Related work.}
The lunar-PNT performance literature is, almost entirely, \emph{user-side}: given a
realized frame and a set of beacons, how accurately can a receiver be located.
\citet{pohlmann2025hybrid} derive error models and Cram\'er--Rao lower bounds for
hybrid satellite/cooperative lunar navigation; open simulators such as LuPNT
\citep{iiyama2023lupnt} target the same receiver-positioning question. These works
take a single, given frame and dynamics as ground truth and do not model the error a
user incurs from provider-to-provider disagreement. On the infrastructure side, the
interoperability standards \citep{lnis2023,ioag2022lunar,israel2020lunanet} and the
reference-frame work define and realize the frames and timescales:
\citet{sosnica2025ilrf} define and realize the International Lunar Reference Frame
(ILRF) and quantify its origin-realization limit---the input, in our companion paper
\citep{baweja2026datum}, to the per-provider frame-error covariance---while
\citet{fienga2024moonlight} survey lunar reference systems, frames, and timescales
in the Moonlight context, noting that ephemeris and frame differences matter for
consistency. Our contribution is distinct from both strands: neither a user-side CRLB
nor a frame realization, but the \emph{cross-provider} budget that sits between
them---how much two providers may differ before the user-side bounds are violated. We
emphasize that the notion of an \emph{irreducible} inter-provider floor is
\emph{derived} here from the structure of the real ephemeris disagreement, not quoted
from any source; where prior surveys remark qualitatively that ephemeris differences
are hard to remove, we make the statement precise, prove which part is removable by a
frame convention and which is not, and measure both.

\paragraph{Contributions.}
We make three contributions, each with an explicit statement of what is proven,
what is independently validated on real data, and what is representative.

\begin{enumerate}
  \item \textbf{A cross-provider consistency-tolerance theorem} (\Proven{};
  Section~\ref{sec:theorem}). We formalize the interop tax as a free-network gauge
  phenomenon: a datum \emph{common} to all providers is unobservable to a
  single-provider user and cancels in provider-versus-provider fusion
  (Lemma~\ref{lem:gauge}), becoming real position error only against an absolute
  external tie; the quantity a fusing or hand-off user actually pays is the
  \emph{differential} datum $\datum_{ab}$ that two providers do not share. That
  differential tax splits---as a measured affine decomposition of the recovered
  rotation---into a \emph{reducible} part removable by a common frame convention and
  an \emph{irreducible} part removable only by a common dynamical reference
  (Lemma~\ref{lem:split}); and a triangle bound on the Helmert point-Jacobian inverts
  into a closed-form consistency tolerance $\tau(B)$ per Helmert parameter for a user
  budget $B$ (Theorem~\ref{thm:tolerance}). The tolerance $\tau(B)$ is exactly the
  number a standard at the AD-5 level must contain.

  \item \textbf{The real inter-ephemeris decomposition} (\Validated{};
  Section~\ref{sec:interephemeris}). We difference the DE440, INPOP21a, and EPM2021
  geocentric-Moon states over 2024--2025 and decompose each pairwise disagreement
  into a Helmert datum and a residual, using the planetary (SSB) positions to
  discriminate a body-common frame-tie from a Moon-specific excess. We find the
  disagreement is an \emph{orientation} effect---a constant rotation explains
  $71$--$94\%$ by amplitude ($91$--$99.7\%$ by variance)---splitting into a reducible
  frame-tie ($\approx\!1$--$3$~nrad) and an irreducible Moon-specific excess
  ($\approx\!5$--$6$~nrad) for the JPL-versus-others pairs. Reported in the
  lever-arm-independent rotations, that excess is a geocentric-Moon-state
  inconsistency worth $\approx\!2$~m at the Earth--Moon lever arm (ephemeris hand-off
  and cislunar users) but only $\approx\!8$~mm for a lunar-surface user. This
  decomposition is \Validated{} against an independent SciPy singular value
  decomposition to a tolerance below $10^{-3}$.

  \item \textbf{An interop budget and a frame-versus-ephemeris design law}
  (\Modelled{}; Section~\ref{sec:budget}). Folding the theorem and the anchor into a
  multi-provider budget, we show that under the frame convention today's standards
  pursue (agree on frame tags and format only) the irreducible rotation is about
  twice the reducible---irreducible fraction
  $\|\rot_{\mathrm{exc}}\|/(\|\rot_{\mathrm{tie}}\|+\|\rot_{\mathrm{exc}}\|)=0.687$---%
  so a common format-and-frame tag leaves the dominant floor: a $\approx\!1.8$~m
  geocentric-Moon-state envelope at the Earth--Moon lever arm, closed only by
  mandating a common (or consistency-bounded) dynamical reference. We report the
  resulting $\tau(B)$ tolerance map for a lunar-surface user.
\end{enumerate}

\paragraph{Honesty scope.}
We separate three kinds of claim. The \emph{theorem}---the gauge lemma, the
reducible/irreducible split as a decomposition of the differential tax, the triangle
bound, and the $\tau(B)$ inversion form---is \Proven{}; the measured \emph{affine}
attribution of the recovered rotation into a frame-tie and a Moon-excess is not (it is
the modelling interpretation flagged below). The \emph{inter-ephemeris
decomposition}---the rotation-only Helmert-difference fit of the real DE440/INPOP21a/%
EPM2021 disagreement---is \Validated{} against an independent numerical oracle. The
\emph{system-level budget}: the multi-provider user model, the total interop budget,
the design-law magnitudes, and the $\tau(B)$ map, are \Modelled{}, because the
``providers'' are ephemerides, a representative analog---no two lunar-PNT providers
fly---and are tagged \modtag{} wherever a system-level number appears. One further
caveat runs throughout and is stated at each occurrence
(Sections~\ref{sec:theorem},~\ref{sec:interephemeris},~\ref{sec:discussion}):
attributing the planet-common rotation to a removable frame-tie is itself a modelling
interpretation; the robust, convention-free claim is the Moon-\emph{excess} rotation,
which no whole-frame choice can remove. Section~\ref{sec:validation} makes the
firewall precise, and every number is regenerated from the committed engine by a
single command (Section~\ref{sec:availability}).
